\documentclass[11pt]{article}
\usepackage[margin=1.1in]{geometry}
\usepackage{amsmath, amssymb, booktabs}

\begin{document}

\begin{table}[htbp]
\centering
\caption{Consumables intensity by career stage: lab consumables spending relative to NIH funding, 2010--2013}
\label{tab:consumables_intensity}
\begin{tabular}{@{}lrrrr@{}}
\toprule
 & Early-career & Late-career & Difference & $p$-value \\
\midrule
\multicolumn{5}{@{}l}{\textit{Panel A. PI level (averages over 2010--2013)}} \\
\addlinespace
Lab consumables spend / NIH funding (\%), mean   & 7.1  & 5.7  & 1.4  & 0.25 \\
\quad median                                      & 7.3  & 4.7  &      &      \\
\quad pooled (sum of spend / sum of NIH)          & 6.2  & 4.8  &      &      \\
High-confidence lab spend / NIH funding (\%), mean & 4.2  & 3.4  & 0.8  & 0.27 \\
\quad pooled                                      & 3.7  & 2.8  &      &      \\
Total FOIA spend / NIH funding (\%), mean         & 12.9 & 9.1  &      &      \\
\addlinespace
Avg.\ annual lab consumables spend (\$)           & 22,799  & 27,136  &      &      \\
Avg.\ annual NIH funding (\$)                     & 365,080 & 569,530 &      &      \\
Median first last-author year                     & 2006    & 1991    &      &      \\
PIs                                               & 25      & 29      &      &      \\
\midrule
\multicolumn{5}{@{}l}{\textit{Panel B. PI-year level}} \\
\addlinespace
Lab consumables spend / NIH funding (\%), mean   & 6.4  & 5.5  & 0.9  & 0.41 \\
\quad median                                      & 5.2  & 4.4  &      &      \\
\quad pooled (sum of spend / sum of NIH)          & 6.0  & 4.7  &      &      \\
High-confidence lab spend / NIH funding (\%), mean & 3.8  & 3.3  & 0.5  & 0.43 \\
\addlinespace
PI-years with positive NIH funding                & 74   & 105  &      &      \\
PI-years with zero NIH funding                    & 20   & 15   &      &      \\
PIs                                               & 28   & 32   &      &      \\
\bottomrule
\end{tabular}
\par\vspace{0.5em}
\begin{minipage}{\textwidth}
\footnotesize
\textit{Notes:} The sample is the FOIA PIs in the main analysis sample
(all journals, rounds 1 and 2) with procurement records in 2010--2013 who are
matched to an NIH RePORTER PI (60 of 77). Lab consumables spend is annual FOIA
procurement spending in categories classified as lab consumables; the
high-confidence measure keeps only categories whose classifier passes the
precision and recall thresholds used in the first stage. NIH funding is the
sum of total costs across the PI's active RePORTER projects in the fiscal
year. Career stage splits PIs at the sample median year of first last-author
publication (2000); early-career PIs are those whose first last-author
publication is after 2000. Panel A averages each PI's spend and NIH funding
over 2010--2013 and reports the ratio of the averages; PIs averaging under
\$10,000 of NIH funding per year are dropped (54 of 60 remain). Panel B uses
PI-years with positive NIH funding. The pooled rows divide the group's total
spend by its total NIH funding. $p$-values are from a two-sided $t$-test of
the early--late difference, clustered by PI in Panel B; Kolmogorov--Smirnov
tests of the distributions give $p = 0.32$ (Panel A) and $p = 0.11$ (Panel B).
Regressing the PI-level intensity on lab age (2014 minus first last-author
year) gives a slope of $-0.11$ percentage points per year (robust SE $0.05$,
$N = 54$). This is a check on consumables intensity, not a budget constraint:
NIH funding understates total research resources, so the ratios are intensity
measures rather than budget shares.
\end{minipage}
\end{table}

\end{document}
